\section{Related work}\label{sec:related}

\emph{Caustic.} Every chord of an irrational orbit of the circular billiard is tangent to the
concentric circle of radius $a$, so the orbit never enters the open disk of radius $a$; this
is the classical caustic of the circular table and is standard \cite{Tab05}.

\emph{Three distances.} The gap structure of $\{n\alpha\}$ is governed by the three-distance
theorem \cite{Sos58,Swier59,Slater67}, with proofs and combinatorial refinements in
\cite{vRav88,AB98} and a recent account of the gap multiplicities in \cite{Leong17}. Blank and
Krikorian proved the analogous three-length theorem for the linear flow on a torus with a slit
\cite{BK93}. Our reduction of the planar tube to the one-dimensional gap law uses only the first
of these.

\emph{Fixed-angle entrance and hitting laws.} For a single irrational rotation, the entrance-time
laws along the denominators $q_k$ need not converge along the full sequence, and Coelho and de
Faria classify the subsequential limits as uniform or piecewise linear \cite{CdF96}. Coelho shows
that tightness of the rescaled hitting-time processes characterizes bounded type
\cite{Coelho04}; Dzhalilov and Khomidov prove piecewise linear limit laws for quadratic
irrational rotations \cite{DK24}; renormalization methods appear in \cite{BI09}. These are
statements about one-dimensional return or hitting times. The planar vacancy area of
Theorem~\ref{thm:B} is a different observable, which we reduce to the one-dimensional gap law.

\emph{Dispersion.} Niederreiter introduced the dispersion
\[\mathrm{Disp}(\{x_n\})=\limsup_{N\to\infty} N d_N,\qquad d_N=\sup_{y\in[0,1]}\min_{n\le N}|y-x_n|,\]
as a measure of denseness \cite{Nie84}. For $x_n=\{n\alpha\}$
it is finite exactly for badly approximable $\alpha$, and Drobot gave a formula for quadratic
irrationals \cite{Dro86}. Niederreiter works on $[0,1]$ rather than on the circle, so the
boundary terms differ from the circular maximal gap used here. Dupain studied the discrepancy of
$\{n\alpha\}$ for the golden ratio \cite{Dup79}. The threshold $c_*$ in
Theorem~\ref{thm:C}(4) is proportional to the circular analogue of this constant, namely
the limsup $\kappa(\beta)$ of $N$ times the maximal circular gap, which equals $1+2/\sqrt5$
for noble angles. Section~\ref{disp:sec} proves that noble angles minimize $\kappa$ and
that $\pi Q\kappa(\beta)$ is the coverage threshold for every irrational angle.

\emph{Lorentz-gas free path.} In the periodic Lorentz gas, the free path in a fixed irrational
direction is controlled by a three-strip partition of the torus minus a scatterer, obtained from
\cite{BK93}; see \cite{MS03,BZ07}. These works compute areas of strips as the scatterer radius
tends to zero, which is closest in method to ours, but the setting is a lattice of scatterers
rather than the chords of one orbit.

\emph{Coverage and billiards.} Fan and Schmeling study coverings of the circle driven by
rotations \cite{FS02}, and Zhao proves almost sure convergence of cover times
\cite{Zhao25}; both concern covering by intervals, not a planar tube. For rectangular billiards,
Don studies the polygons cut out by a truncated orbit and the number of distinct areas
\cite{Don12}. Hitting statistics for integrable flows are treated in \cite{DMS17} for
absolutely continuous families of frequencies, and do not apply to one fixed angle.

\emph{What is new.} We did not find in these sources the exact two-branch description of the
$w$-neighborhood of $N$ complete chords, the resulting formula for the uncovered area, or its
extremal behavior at Fibonacci $N$ for the golden angle. To our knowledge these are new;
the one-dimensional ingredients are classical.

% BIBITEMS
% \bibitem{BK93} S.~Blank and N.~Krikorian, Thom's problem on irrational flows, \emph{Internat. J. Math.} \textbf{4} (1993), 721--726.
% \bibitem{Dro86} V.~Drobot, On dispersion and Markov constants, \emph{Acta Math. Hungar.} \textbf{47} (1986), 89--93.
% \bibitem{Dup79} Y.~Dupain, Discr\'epance de la suite $(\{n\alpha\})$, $\alpha=(1+\sqrt5)/2$, \emph{Ann. Inst. Fourier (Grenoble)} \textbf{29} (1979), no.~1, 81--106.
% \bibitem{Don12} H.~Don, Polygons in billiard orbits, \emph{J. Number Theory} \textbf{132} (2012), no.~6, 1151--1163; arXiv:1106.2030.
% \bibitem{FS02} A.-H.~Fan and J.~Schmeling, Coverings of the circle driven by rotations, in: \emph{Dynamical Systems from Number Theory to Probability II} (A.~Khrennikov, ed.), V\"axj\"o University Press, 2002, 7--15.
% \bibitem{Nie84} H.~Niederreiter, On a measure of denseness for sequences, in: \emph{Topics in Classical Number Theory, Vol.~II (Budapest, 1981)}, Colloq. Math. Soc. J\'anos Bolyai \textbf{34}, North-Holland, Amsterdam, 1984, 1163--1208.
% \bibitem{Tab05} S.~Tabachnikov, \emph{Geometry and Billiards}, Student Mathematical Library \textbf{30}, Amer. Math. Soc., Providence, RI, 2005.
